\section{Preliminaries}
\label{sec:pre}
We first start by formally defining $2$-CSP.
\begin{definition}[2CSP]\label{def:2csp}
An instance of arity-2 constraint satisfaction problem (2CSP) is a tuple $\Pi=(X,\Sigma,\Phi)$, where:
\begin{itemize}
\item $X=\{x_1,\ldots,x_k\}$ is the set of variables;
\item $\Sigma=\{\Sigma_{x_1},\ldots,\Sigma_{x_k}\}$ is the set of their respective domains of values. We use $|\Sigma|$ to denote the maximum size of any domain, and call it the alphabet size of $\Pi$. 
\item $\Phi=\{\phi_1,\ldots,\phi_m\}$ is the set of constraints. Each constraint $\phi_i$ is 
 ordered tuple $\langle w_i,R_i\rangle$, where $w_i=(w_{i,1},w_{i,2}) \in X^2$ is a pair of variables, and $R_i$ is a 2-ary relation on $\Sigma_{w_{i,1}}$ and $\Sigma_{w_{i,2}}$. 
\end{itemize}

An assignment of the 2CSP instance is a function from all variables in $X$ to their respective domains. A constraint $\phi_i$ is said to be satisfied by an assignment $\sigma$ if $(\sigma(w_{i,1}),\sigma(w_{i,2}))\in R_i$. We say an assignment $\sigma$ satisfies the 2CSP instance if all constraints are satisfied by $\sigma$. 
\end{definition}

For each constraint $\phi_i$, we can without loss of generality assume $w_{i,1} \neq w_{i,2}$, since unary constraints can be removed by restricting the domain of that variable.

We define $r$-list satisfiability, which generalizes the above satisfiability. 
\begin{definition}[$r$-List Satisfiability]
    Given a 2CSP instance $\Pi=(X,\Sigma,\Phi)$, a multi-assignment is a function mapping each variable to a subset of its domain. We define the size of a multi-assignment $\sigma$ as $\max_{x \in X} |\sigma(x)|$.
    
    We say a multi-assignment $\sigma$ $r$-list satisfies $\Pi$ if $\sigma$ is of size at most $r$, and for every constraint $\phi_i$, there exists a pair of values $u \in \sigma(w_{i,1})$ and $v \in \sigma(w_{i,2})$, such that $(u,v) \in R_i$.
\end{definition}

Normal satisfiability can be viewed as $1$-list satisfiability. Note that as $r$ increases, it becomes easier to $r$-list satisfy a 2CSP instance. 

The relations in the $2$CSPs that we construct using the direct product (see \Cref{def:prod}) satisfy a useful structural property, namely, rectangular relations. 
\begin{definition}[Rectangular Relation]\label{def:rect}
    A relation $R \subseteq A \times B$ is said to be rectangular if there is a set $C$ and functions $\pi: A \to C$ and $\sigma: B \to C$ such that $(a,b) \in R$ if and only if $\pi(a)=\sigma(b)$. Equivalently, $R$ is rectangular if for all $a,a' \in A$ and $b,b' \in B$ such that $(a,b)\in R, (a,b')\in R$, and $(a',b) \in R$, we have $(a',b')\in R$.
\end{definition}

Rectangular relations can be informally viewed as consistency checks, and they are often satisfied by $2$CSPs in product constructions. Projection relation, a stronger version of rectangular relation, is ubiquitous in PCP-based hardness reductions.

We now formally define the direct product construction that we use in our proof. 
\begin{definition}[Direct Product Construction]\label{def:prod}
    Given a 2CSP instance $\Pi=(X,\Sigma,\Phi)$, its $t$-wise direct product, denoted as $\Pi^{\odot t}$, is the following 2CSP instance $(X',\Sigma',\Phi')$:
    \begin{itemize}
        \item $X'=\binom{X}{t}$, where each variable is a $t$-sized subset of variables in $\Pi$. 
        \item The domain of each variable $S \in X'$ is the set of all partial satisfying assignments for $S$ in $\Pi$, i.e., all function $\sigma$ that maps each $x \in S$ to its domain in $\Pi$, such that all constraints in $\Pi$ induced by $S$ are satisfied.
        \item $\Phi'$ has of a consistency constraint for each pair of distinct variables in $X'$. For $S,T \in X'$, the assignments $\sigma_S,\sigma_T$  satisfy the constraint if and only if they are consistent on the values for variables in $S \cap T$.
    \end{itemize}
\end{definition}

Our results are based on the hypothesis $\wone \neq \fpt$, which is closely related to  \kclique{}, a fundamental problem in parameterized complexity theory. 
\begin{definition}[\kclique{}]
    An instance of (multicolored) \kclique{} problem is an undirected graph $G=(V=V_1\dot\cup\ldots\dot\cup V_k,E)$, where each $V_i$ is an independent set. The goal is to decide whether we can find $v_1 \in V_1,\ldots,v_k \in V_k$ which form a clique. For $r>1$, the $r$-gap version of \kclique{} asks to distinguish between the following two cases:
    \begin{itemize}
        \item (Yes) There exists $v_1 \in V_1,\ldots,v_k \in V_k$ which form a clique.
        \item (No) The maximum clique in $G$ has size at most $k/r$.
    \end{itemize}
\end{definition}

The $\wone\neq\fpt$ hypothesis states that for any computable function $f$, no algorithm can solve a $\kclique{}$ instance with size $n$ in $f(k)\cdot n^{O(1)}$ time. For a $\kclique{}$ instance $G=(V=V_1\dot\cup\ldots\dot\cup V_k,E)$, we can build $k$ variables $x_1,\ldots,x_k$, letting $V_i$ be the domain of variable $x_i$ and making the edge set between $V_i$ and $V_j$ the constraint relation for $x_i$ and $x_j$. Thus,  $G$ corresponds to a 2CSP instance with $k$ variables and alphabet size at most $n$. It is easy to see that there is a clique of size $k$ in $G$ if and only if the 2CSP instance is satisfiable. Thus, we can restate the $\wone\neq\fpt$ hypothesis as follows. 

\begin{hypothesis}[$\wone\neq\fpt$]
    For any computable function $f$, no algorithm can decide whether a given 2CSP instance $\Pi=(X,\Sigma,\Phi)$ is satisfiable in $f(|X|)\cdot |\Sigma|^{O(1)}$ time.
\end{hypothesis}

The approximation version (with respect to the fraction of constraints that can be satisfiable) of $\wone\neq\fpt$ is called Parameterized Inapproximability Hypothesis (PIH).
\begin{hypothesis}[Parameterized Inapproximability Hypothesis (PIH)\cite{LRSZ20}] \footnote{Note that the original PIH in \cite{LRSZ20} states that constant approximating 2CSP parameterized by $|X|$ is $\wone$-hard. Here we use a relaxed form.}
    For any computable function $f$ and some constant $\varepsilon>0$, no algorithm can given as input a 2CSP instance $\Pi=(X,\Sigma,\Phi)$, distinguish between the following two cases in $f(|X|)\cdot |\Sigma|^{O(1)}$ time:
    \begin{itemize}
        \item (Completeness) $\Pi$ is satisfiable.
        \item (Soundness) Any assignment of $\Pi$ violates at least $\varepsilon$ fraction of constraints.
    \end{itemize}
\end{hypothesis}

We formally define \ksetcover{}, the parameterized version of the classical {\sc SetCover} problem, and the exact version of it. 
\begin{definition}[\ksetcover{}, \kexact{}]
    An instance of \ksetcover{} problem is a tuple $\Pi=(\mathcal S,U)$, where $\mathcal S$ is a collection of subsets $\{S_1,\ldots,S_n\}$ over the universe $U$, and the goal is to decide whether there are $k$ sets in $\mathcal S$, whose union is $U$. For $r>1$, the $r$-gap version of \ksetcover{} asks to distinguish between the following two cases:
    \begin{itemize}
        \item (Yes) There are $k$ sets whose union is $U$.
        \item (No) The union of any $r\cdot k$ sets is not $U$.
    \end{itemize}
    Furthermore, if in the yes case, the $k$ sets are non-intersecting, i.e., they form a partition of $U$, then we also denote this gap problem as \kexact{}.
\end{definition}

Finally, we define the $(T,m)$-set gadget, an important gadget used in reductions to \ksetcover{} and \kexact{}.

\begin{definition}[$(T,m)$-Set Gadget]
    A $(T,m)$-set gadget consists of a universe $\mathcal M$ and some of its subsets $C_1,\ldots,C_{m}$ with the following property: Every collection of at most $T$ sets out of $C_1,\overline {C_1},C_2,\overline{C_2},\ldots,C_m,\overline{C_m}$ that is a set cover for $\mathcal M$ must include both $C_i$ and $\overline{C_i}$ for some $i$.
    %The $(\kappa,\rho)$-hypercube partition system consists of the universe $\mathcal M$ and a collection of subsets $\{P_{x,y}\}_{x \in [\rho],y \in [\kappa]}$, where $\mathcal M=[\kappa]^\rho$ and $P_{x,y}=\{z \in \mathcal M:z_x=y\}$.
\end{definition}

It was proved in \cite{LY94} that a $(T,m)$-set gadget can be constructed efficiently:

\begin{lemma}
\label{lem:construct-set-gadget}
    There is an algorithm that given any $T,m$, runs in time $\text{poly}(m,2^T)$ and outputs a $(T,m)$-set gadget with universe size $\text{poly}(m,2^T)$.
\end{lemma}

%The universe $\mathcal M$ consists of all functions from $[\rho]$ to $[\kappa]$, and is of size $\kappa^{\rho}$. Each subset $P_{x,y} (x \in [\rho],y \in [\kappa])$ consists of all functions mapping $x$ to $y$. It can be observed that one can cover the universe by picking all $\kappa$ subsets from some row $x \in [\rho]$, and this is the only way to cover the universe. In other words, even if we include $\kappa-1$ subsets from every row, it is not possible to cover the universe.
